\documentclass[11pt]{article}
\usepackage{cse216,hyperref}
\def\title{HW 3}
\begin{document}
\maketitle

\section*{Due October 20 at 3:00 pm}

\textbf{(4 questions, 185 points total)}

\begin{qunlist}

\q{20}{Resolution}
\begin{qparts}
\item
Which of the following clauses can be resolved with each other, and what are the resulting clauses?
\begin{enumerate}[(A)]
\item $x \lor y \lor z$
\item $y$
\item $\lnot y$
\item $\lnot y \lor z$
\end{enumerate}

\begin{answer}

\end{answer}

\item
Give a resolution refutation for the UNSAT formula with the following clauses:
\begin{align*}
c_1=\;&\lnot x \lor \lnot y \lor z \\
c_2=\;&\lnot x \lor w \\
c_3=\;&x \lor \lnot y \\
c_4=\;&\lnot y \lor \lnot z \lor \lnot w \\
c_5=\;&x \lor y \\
c_6=\;&y \lor \lnot w
\end{align*}
You may write  $c_i \odot c_j$ to denote the resolution of two clauses.

\begin{answer}

\end{answer}
\end{qparts}


\q{60}{CDCL}\\
Consider running CDCL on the formula in (1b) above, with a variable selection heuristic that prefers $z > w > x > y$ and a phase selection heuristic that always picks the negative polarity.
\begin{qparts}
\item
Up to the first conflict, what decisions are made?

\begin{answer}

\end{answer}

\item
Construct the implication graph for the first conflict.

\begin{answer}

\end{answer}

\item
What are the possible conflict clauses that could be learned from this conflict? Which of these is asserting if we backtrack one level?

\begin{answer}

\end{answer}

\item
If we use the first UIP heuristic, which conflict clause will CDCL learn, and what decision level will it backtrack to?

\begin{answer}

\end{answer}

\item
Construct the implication graph for the next conflict.

\begin{answer}

\end{answer}

\item
CDCL terminates after these two conflicts. Why? How can you use the implication graphs from each conflict to build a resolution refutation of the original formula?

\begin{answer}

\end{answer}

\end{qparts}

\q{70}{First-order logic}\\
For each of the following formulas, state whether they are satisfiable or valid (or neither).
If a formula $\phi$ is satisfiable, give a model, i.e.~a structure $M$ such that $M \models \phi$.
If $\phi$ is not valid, give a structure $M$ such that $M \not \models \phi$.
For each formula, assume the signature of the logic --- the constant, function, and relation symbols --- consists just of those symbols used in the formula.
We use $c$, $f$, and $R$ as constant, function, and relation symbols respectively.
\begin{qparts}
\item
$\forall x . \forall y . (f(x, y) = f(y, x))$

\begin{answer}

\end{answer}

\item
$\exists x. (x = x)$ \quad \emph{(hint: can you pick the universe of the structure so that this is false?)}

\begin{answer}

\end{answer}

\item
$\forall x. \forall y. (x = y)$

\begin{answer}

\end{answer}

\item
$\exists x . (x = x) \land \forall x. \forall y. \lnot (x = y)$

\begin{answer}

\end{answer}

\item
$\forall x . \exists y . (f(x) = y)$

\begin{answer}

\end{answer}

\item
$\exists x. (R(x) \land \lnot R(f(x)))$

\begin{answer}

\end{answer}

\item
$\left[ R(c) \land \forall x . (R(x) \rightarrow R(f(x))) \right] \rightarrow R(f(f(f(c))))$

\begin{answer}

\end{answer}
\end{qparts}

\q{35}{Arrays}\\
Consider the theory of bitvectors (say, 32-bit) and arrays whose indices and values are bitvectors (sometimes called ABV).
This theory has two kinds of objects, or \emph{sorts}: bitvectors and arrays.
Let's use the notation $\forall x : Array .$ and $\forall x : BV .$ to quantify over arrays and bitvectors respectively.
So for example if we use the C notation $>>$ for the right shift operator, then $\forall x : BV . (x >> 32 = 0)$ would be a valid formula asserting that for all 32-bit bitvectors, right shifting by 32 bits yields the zero bitvector.
\begin{qparts}
\item
Write an SMT formula $\phi(x)$ which is true if and only if the array $x$ is sorted in increasing order.

\begin{answer}

\end{answer}

\item
Write a formula $\phi(x)$ which asserts that every element of the array $x$ has the same value.

\begin{answer}

\end{answer}

\item
Write a formula with no free variables (a \emph{sentence}) asserting that if element 0 of an array is even, and for every $i$ if element $i$ is even then element $i+1$ is even, then all elements of the array are even.
Is this formula valid?

\begin{answer}

\end{answer}

\end{qparts}

\q{optional, bonus problem; 0}{GCHQ puzzle}\\
The British signals intelligence organization GCHQ released a \href{https://web.archive.org/web/20160330032843/http://www.gchq.gov.uk/press_and_media/news_and_features/Pages/Directors-Christmas-puzzle-2015.aspx}{puzzle} in December 2015, the correct solution of which led to a sequence of increasingly difficult puzzles (presumably with your employment as a codebreaker at the end --- alas, the puzzles have been taken down).
They described it as ``a brain-teasing puzzle that seems certain to exercise the grey matter of participants over the holiday season''.
Don't exert yourself: solve it with SAT.

\end{qunlist}
\end{document}